\documentclass[11pt]{article}

\usepackage[margin=25mm]{geometry}

\usepackage{amsmath,amssymb,fontspec,unicode-math}

\IfFontExistsTF{DejaVu Serif}{\setmainfont{DejaVu Serif}}{\setmainfont{Latin Modern Roman}}

\IfFontExistsTF{Latin Modern Math}{\setmathfont{Latin Modern Math}}{\setmathfont{DejaVu Math TeX Gyre}}

\usepackage{xurl,hyperref}

\setlength{\parindent}{0pt}\setlength{\parskip}{6pt}

\begin{document}

\title{RAST: Reasoning Activation in LLMs via Small-model Transfer}

\author{ELVIA preliminary AI analysis}\date{21 October 2026}\maketitle

Preliminary assessment: Minor.\par

Independent scientific review remains pending. Author adoption has not been requested.

The teacher-forced PCR study does not by itself establish the proposed distributional/knowledge interpretation, but that story is framed as a hypothesis rather than a proved bound. Local analysis-script, memory arithmetic and reporting corrections do not demonstrate failure of the central measured transfer results. The strict assessment is Minor.

\section*{Review scope}

Full original concerns, fresh current proceedings methods/results and implicated appendices, plus current pinned analysis, evaluation, prompt and single-expert decoding source inspected read-only. No author or historical verifier program was executed.

\section*{Problem 1: Teacher-forced PCR does not establish small distributional changes}

Source: PDF pp.2–3 Eq 1–2; Appendix C.1 p.26

Assessment: supported; impact: minor

\subsection*{Source claim (paraphrase)}

The preliminary PCR study motivates the hypothesis that RL changes probabilities mainly at a small set of reasoning-related tokens.

\subsection*{Analysis}

Take a vocabulary of \(100\) tokens. Let two next-token laws be \[p=(0.02,0.98/99,\ldots,0.98/99)\] and \[q=(0.99,0.01/99,\ldots,0.01/99)\]. Both have the same unique argmax, but \(\|p-q\|_{\rm TV}=0.97\). Repeating this construction on every teacher-forced prefix gives perfect argmax preservation while distributional change remains large. Teacher-forced prefixes also differ from prefixes generated by a model's own earlier tokens.

\subsection*{Finding and scope}

Identical teacher-forced argmax decisions do not establish small probability changes or free-running path agreement. The paper presents a motivating hypothesis, not a formal total-variation or knowledge-preservation theorem; the counterexample limits that interpretation without refuting measured reasoning gains.

\subsection*{Proposed correction}

Describe conditional top \(1\) agreement; measure distributions/likelihood and perform controlled token-position interventions.

Proposed correction; independent verification and author adoption remain unestablished.

\section*{Problem 2: Raw logit-difference cosine is not invariant to probability-preserving shifts}

Source: PDF p.7 Eq 6/Fig 4; analysis/similarity.py:89–105

Assessment: needs\_evidence; impact: needs\_evidence

\subsection*{Source claim (paraphrase)}

Raw logit-difference cosine identifies probability-distribution alignment.

\subsection*{Analysis}

Let centered logit differences be \(a=(1,-1,0)\) and \(b=(1,0,-1)\), with cosine \(1/2\). Adding a common scalar to all logits preserves each softmax law, yet large same-sign or opposite-sign constant shifts make the raw difference cosine approach \(1\) or \(-1\). Thus this cosine is not determined by the probability laws.

\subsection*{Finding and scope}

This is a scope/provenance gap, not proof that the stronger empirical conclusion is false.

\subsection*{Proposed correction}

Center logit differences or use explicit probability/log-ratio metrics; align semantic vocabulary and report model/trajectory uncertainty.

Proposed correction; independent verification and author adoption remain unestablished.

\section*{Problem 3: Cosine analysis cumulatively duplicates answer prefixes}

Source: analysis/similarity.py:76–114

Assessment: supported; impact: minor

\subsection*{Source claim (paraphrase)}

The cosine-analysis loop advances prefixes one token at a time.

\subsection*{Analysis}

The author similarity loop appends the accumulating answer prefix repeatedly. For an initial context of length \(2\), it produces prefix lengths \(2,3,5,8\), rather than \(2,3,4,5\). The evaluated sequences therefore repeat earlier answer tokens; this is a local analysis-script defect, not an estimate of corrected paper conclusions.

\subsection*{Finding and scope}

The bounded local discrepancy or counterexample is supported; no change in published benchmark scores is quantified.

\subsection*{Proposed correction}

Keep prompt\_ids immutable and append answer[: i] to it, or increment cached prefix by one token.

Proposed correction; independent verification and author adoption remain unestablished.

\section*{Problem 4: PCR coefficient shadowing, denominator mismatch and one-run averaging}

Source: analysis/greedy\_delta.py:49–80; simplelr\_math\_eval/math\_eval\_pred.py:637–646

Assessment: supported; impact: minor

\subsection*{Source claim (paraphrase)}

PCR uses the supplied coefficient and correct position/run denominators.

\subsection*{Analysis}

The coefficient argument named i is overwritten by a loop index of the same name. In addition, a sequence with \(99\) perfectly preserved checked predictions is divided by \(100\), and the averaging branch runs only range(\(1\)), yielding a single value with zero reported spread. These contracts do not implement the intended coefficient, denominator and repeated-run summary.

\subsection*{Finding and scope}

The bounded local discrepancy or counterexample is supported; no change in published benchmark scores is quantified.

\subsection*{Proposed correction}

Separate lambda and token index, normalize valid positions, explicitly load \(32\) unique aligned run files and publish exact producing scripts.

Proposed correction; independent verification and author adoption remain unestablished.

\section*{Problem 5: Minerva recovery rate arithmetic error}

Source: PDF p.5 Table 1; p.4 Eq 5

Assessment: supported; impact: minor

\subsection*{Source claim (paraphrase)}

The Minerva recovery rate follows Eq.5.

\subsection*{Analysis}

The Minerva recovery fraction from the displayed entries is \((25.1-21)/(33.6-21)\simeq0.3253968\), not \(0.841\). The other \(35\) displayed recovery entries and Table 4 averages do not show this arithmetic discrepancy.

\subsection*{Finding and scope}

The bounded local discrepancy or counterexample is supported; no change in published benchmark scores is quantified.

\subsection*{Proposed correction}

Correct to \(32.5\%\) or regenerate from raw unrounded values and reconcile source numbers.

Proposed correction; independent verification and author adoption remain unestablished.

\section*{Problem 6: Decoding and prompt controls confound recovery interpretation}

Source: PDF p.5 Table 1 caption/footnote; p.25 Table 5/Appendix A.2; vllm/tests/RAST\_decode/run.py:54–60

Assessment: needs\_evidence; impact: needs\_evidence

\subsection*{Source claim (paraphrase)}

The decoding/prompt controls isolate the claimed recovery effect.

\subsection*{Analysis}

Table 1 and Appendix A disclose greedy base decoding versus temperature-\(1\), top-\(p=0.95\) RAST/RL sampling. Thus the recovery ratio compares complete systems and does not isolate the logit adjustment under a common decoding policy. Figure 6 gives approximately \(70.5\) on MATH500 at \(\lambda=0\) with sampling, versus the table’s greedy base \(68.6\); this is illustrative, not an exact decomposition of every gain. The dagger note and Appendix A.2 give differing descriptions of the small-model prompt, while the current run.py fixes use\_chat to false. The producing run configuration is required to settle prompt attribution.

\subsection*{Finding and scope}

This is a scope/provenance gap, not proof that the stronger empirical conclusion is false.

\subsection*{Proposed correction}

Report all-greedy and matched-sampling panels; publish exact input strings and decoding configs per model.

Proposed correction; independent verification and author adoption remain unestablished.

\section*{Problem 7: GRPO resource model includes critic and inconsistent CPU arithmetic}

Source: PDF p.29 AppendixD; p.8 Table 4; original DeepSeekMath §4.1.1

Assessment: supported; impact: minor

\subsection*{Source claim (paraphrase)}

The resource model correctly accounts for GRPO components and optimizer memory.

\subsection*{Analysis}

GRPO does not require a learned critic. The stated CPU-memory arithmetic for two FP32 Adam state arrays, \(14\times10^9\) parameters and three models is \(2\cdot14\times10^9\cdot4\cdot3=336\times10^9\) bytes, not \(156\) GB. Counting optimizer states for frozen models also needs justification. This corrects the resource model rather than measuring the actual runtime/memory of any run.

\subsection*{Finding and scope}

The bounded local discrepancy or counterexample is supported; no change in published benchmark scores is quantified.

\subsection*{Proposed correction}

Profile actual active copies/trainable states by phase, separate CPU/GPU, rollout/reference/old-policy and activations, and recompute savings.

Proposed correction; independent verification and author adoption remain unestablished.

\section*{Problem 8: Throughput reduction confused with latency increase}

Source: PDF p.8 Table 4 and inference-speed paragraph

Assessment: supported; impact: minor

\subsection*{Source claim (paraphrase)}

A throughput reduction equals the stated latency increase.

\subsection*{Analysis}

For equal generated-token counts and comparable hardware/batching, generation time is inversely proportional to throughput. The current Table 4 rates give RAST time increases \(243.1/200.3-1\approx21.37\%\), \(243.1/202.7-1\approx19.93\%\), and \(249.5/216.1-1\approx15.46\%\). The corresponding throughput reductions are \(17.61\%\), \(16.62\%\), and \(13.39\%\). These do not yield the stated \(13\)–\(16\%\) latency range. This is an equal-work arithmetic correction, not an independent measurement of request latency; actual output lengths, queueing and prompt processing also matter.

\subsection*{Finding and scope}

The bounded local discrepancy or counterexample is supported; no change in published benchmark scores is quantified.

\subsection*{Proposed correction}

Report throughput as throughput or use inverse-rate conversion for equal workload; measure end-to-end latency on matched requests.

Proposed correction; independent verification and author adoption remain unestablished.

\section*{Problem 9: Improved worked example reverses exponential inequalities}

Source: PDF p.29 Figure 10

Assessment: supported; impact: minor

\subsection*{Source claim (paraphrase)}

The displayed improved proof uses correct exponential comparisons.

\subsection*{Analysis}

In the worked exponential comparison, \(x=0\) gives \(9^x=1>10^x-19=-18\), while \(x=3\) gives \(729<981\). The generated explanation reverses those inequalities. The root \(x=2\) is correct; the proof narrative is not.

\subsection*{Finding and scope}

The bounded local discrepancy or counterexample is supported; no change in published benchmark scores is quantified.

\subsection*{Proposed correction}

Keep the actual erroneous output visible and correct its attribution. A valid root argument compares increasing \((10/9)^x\) with decreasing \(1+19\cdot9^{-x}\).

Proposed correction; independent verification and author adoption remain unestablished.

\section*{Problem 10: KL \(837.9\) lacks units and producing-code provenance}

Source: PDF p.28 Appendix C.5; analysis/similarity.py:139–141 commented code

Assessment: needs\_evidence; impact: needs\_evidence

\subsection*{Source claim (paraphrase)}

The reported KL value has a defined producing formula and units.

\subsection*{Analysis}

The available author analysis snippet is commented out and computes a scaled quantity \(\operatorname{KL}(\mathrm{RAST}\|\mathrm{base})\,10^9/V\). If this were the rule used with \(V=152064\), a displayed \(837.9\) would correspond to \(0.1274144256\) nats. The original plot-producing code is not established, so this is a units/provenance question rather than a replacement value for the figure.

\subsection*{Finding and scope}

This is a scope/provenance gap, not proof that the stronger empirical conclusion is false.

\subsection*{Proposed correction}

Release logits/logprobs, vocabulary mask, direction, scaling, exact script and unscaled KL.

Proposed correction; independent verification and author adoption remain unestablished.

\section*{Problem 11: Monotonic majority and uniformly increasing token-frequency statements conflict with own plots}

Source: PDF p.26 Appendix C.3; p.27 Figure 8; p.8 prose/p.9 Figure 5

Assessment: supported; impact: minor

\subsection*{Source claim (paraphrase)}

The stated majority and token-frequency trends are monotone.

\subsection*{Analysis}

For independent binary answers with correctness probability \(p=0.4\), majority-vote correctness at \(1,3,5\) draws is \(0.4,0.352,0.31744\). Majority voting is therefore not universally monotonic in correctness. The paper's own token-frequency plots also include decreases, such as MATH branching \(0.87\to0.46\) and AIME self-verification \(3.20\to2.30\).

\subsection*{Finding and scope}

The bounded local discrepancy or counterexample is supported; no change in published benchmark scores is quantified.

\subsection*{Proposed correction}

Describe nonmonotonic observed trends and movement toward RL frequencies; correct AIME subplot label.

Proposed correction; independent verification and author adoption remain unestablished.

\section*{Problem 12: Uncertainty, lexical proxy, transfer scope and benchmark versions need stronger specification}

Source: PDF §§3–4, Table 8; reasoning\_tokens.py; Appendix A.1

Assessment: needs\_evidence; impact: needs\_evidence

\subsection*{Source claim (paraphrase)}

The reported uncertainty and proxies establish the claimed general transfer.

\subsection*{Analysis}

The repeated generations share benchmark questions, so \(32\) runs are not \(32\) independent benchmark samples. The paper does show standard-deviation information in Figure 6; a claim of no uncertainty reporting anywhere would be false. The author reasoning\_tokens.py averages keyword occurrences per response, rather than per output token. Different response lengths can therefore contribute to the reported lexical differences, and broad words such as “but” are imperfect behavioral proxies. The experiments transfer across scales within Qwen and within Llama; they do not establish arbitrary transfer between incompatible tokenizers. LiveCodeBench is described by a count of \(880\) questions without a fully frozen release/window. These are limits on statistical, mechanism and version attribution, not demonstrated data contamination.

\subsection*{Finding and scope}

This is a scope/provenance gap, not proof that the stronger empirical conclusion is false.

\subsection*{Proposed correction}

Publish paired prompt-level uncertainty, length-adjusted/semantic behavior measures, versioned benchmark IDs and explicit tokenizer assumptions.

Proposed correction; independent verification and author adoption remain unestablished.

\section*{Problem 13: Original MATH split counts incorrect}

Source: PDF p.24 Appendix A.1; original MATH paper§3, p.4

Assessment: supported; impact: minor

\subsection*{Source claim (paraphrase)}

The original MATH split and related benchmark labels are correctly described.

\subsection*{Analysis}

Original MATH is \(7500\) train/5000 test, not \(8000\)/4500. MBPP+ paragraph also misnames HumanEval+.

\subsection*{Finding and scope}

The bounded local discrepancy or counterexample is supported; no change in published benchmark scores is quantified.

\subsection*{Proposed correction}

Correct original split and dataset name; separately identify the \(500\)-problem evaluation subset.

Proposed correction; independent verification and author adoption remain unestablished.

\section*{Limits}

Independent scientific adjudication remains pending.\par

No author code, model training or full benchmark evaluation was executed in this revision.\par

Version-specific source defects do not establish which code produced the published experiments.\par

Corrected benchmark scores, empirical effect sizes and author adoption remain unverified unless explicitly stated.\par

No new author-code execution, training, empirical replication or contact.\par

Historical verification outputs were inspected but not rerun.\par

Missing run provenance and unresolved empirical interpretations remain explicit.\par

\section*{Source identity}

NeurIPS 2025 current proceedings PDF (checked 21 October 2026)

\url{https://proceedings.neurips.cc/paper_files/paper/2025/file/4757e1d1a5890c42c85e9864daf12989-Paper-Conference.pdf}

Reviewed PDF SHA-256: \nolinkurl{9d6d6bc032c0ed42988afb762ad506d9d11d84e57641af064120eae60995e358}

This current primary is the same research work as the predecessor, identified by matching title, authors and technical content, on the stated current publication channel. Its exact inspected bytes are identified by SHA-256; the earlier immutable source identity remains in the predecessor record. The download timestamp is recorded from this fresh primary-source retrieval.

\end{document}
